\documentclass[12pt]{article}
\usepackage[a4paper, margin=2cm]{geometry}
\usepackage[T1]{fontenc}
\usepackage[czech]{babel}
\usepackage[utf8]{inputenc}
\usepackage{csquotes}
\usepackage{xcolor}
\usepackage{tikz}
\usepackage{paralist}
\usepackage{scrextend}
\usepackage{newtxtext} % Times New Roman
\usepackage{tocloft} % table of contents
\usepackage{listings} % code listings

\definecolor{codegreen}{rgb}{0,0.6,0}
\definecolor{codegray}{rgb}{0.5,0.5,0.5}
\definecolor{codepurple}{rgb}{0.58,0,0.82}
\definecolor{backcolour}{rgb}{0.95,0.95,0.92}

\lstdefinestyle{mystyle}{
	backgroundcolor=\color{backcolour},
	commentstyle=\color{codegreen},
	keywordstyle=\color{magenta},
	numberstyle=\tiny\color{codegray},
	stringstyle=\color{codepurple},
	basicstyle=\ttfamily\footnotesize,
	breakatwhitespace=false,
	breaklines=true,
	captionpos=b,
	keepspaces=true,
	numbers=left,
	numbersep=5pt,
	showspaces=false,
	showstringspaces=false,
	showtabs=false,
	tabsize=2
}

\lstset{style=mystyle}

\let\e\expandafter

%TODO{spellcheck}

\def\showtodo{\immediate\write16{}}
\def\TODO#1{%
	\textbf{ [todo: #1] }%
	\e\xdef\e\showtodo\e{\showtodo\immediate\write16{todo on line \number\inputlineno: #1}}%
}

\setlength{\parindent}{0pt}
\setlength{\parskip}{\baselineskip}%

\def\file#1{\texttt{#1}}
\def\code#1{\texttt{#1}} %\TODO{use a~verbatim environment}

\def\title{Název práce}
\def\author{Michael Jarvis}
\def\year{2025/2026}
\def\class{8.V}

\begin{document}
\begin{titlepage}
	\vspace*{\fill}
	\centering
	{\Huge\title\par}
	\vspace{3\baselineskip}
	{\Large Dokumentace maturitní práce\par}
	\vspace{2\baselineskip}
	{\Large\author\par}
	\vfill\vfill
	{\large školní rok: \year}\hfil{\large třída: \class}
	\vspace*{3cm}
\end{titlepage}

\newpage \TODO{this page looks bad}
{\Large Zadání maturitní práce:\par}

Kompiler\\
Cílem práce je vytvořit kompiler pro vlastní jazyk podobný C~za použití knihoven flex, bison a~LLVM toolchain. Jazyk by měl podporovat standardní operace s~celými čísly (sčítání, odčítání, násobení, dělení, …), definice a volání funkcí, cykly, strukturované datové typy (“struct” v jazyku C).
Projekt bude napsaný v~jazyce C~a~bude pro operační systém Linux.

Teoretická část - Stavové automaty v~parserech


\newpage
\vspace*{\fill}
\begin{center}
Prohlašuji, že jsem na práci pracoval samostatně pouze za pomoci použitých zdrojů a~že
v~práci i~v~dokumentaci jasně vymezuji, které části kódů jsou mým originálním dílem, které
jsou upravenou verzí a které jsou převzaty v plném rozsahu.
\end{center}
\vspace{2cm}
\begin{flushright}
\_\_\_\_\_\_\_\_\_\_\_\_\_\_\_\_\_\_\_\_\_\_\_\_\_\_\_\_\_\_\_\_
\vspace{5pt}\\
podpis žáka \hspace*{4cm}
\end{flushright}
\vspace*{\fill}

\newpage
\renewcommand\contentsname{Obsah}
\renewcommand\cftsecaftersnum{.}
\renewcommand\cftsecleader{\cftdotfill{\cftdotsep}}
\tableofcontents






\newpage
\section{TEORETICKÁ ČÁST} \TODO{teoreticka cast}
\textit{zpracována podle zadání}

\section{CÍLE PRÁCE}
Cílem této práce bylo vytvořit program, který jako vstup dostane soubor se zdrojovým kódem a~přeloží ho do kódu LLVM, který pomocí knihovny LLVM toolchain přeloží do strojového kódu.
V~případě chyby ve vstupním zdrojovém kódu vypíše chybovou hlášku. Také by měl umět v~některých případech vypisovat pomocná varování.

Jazyk, který program zpracovává je velmi podobný jazyku C, ale je mnohem omezenější.
Měl by umět základní operace s~čísly a~proměnnými (včetně globálních proměnných), funkce a~strukturované datové typy.


\section{ZPŮSOBY ŘEŠENÍ A POUŽITÉ POSTUPY} \TODO{}
%\textit{JAK JE TO ŘEŠENO}\\
%\textit{podrobný popis programu, použité algoritmy, datové struktury, způsoby testování funkčnosti, ošetření chyb}

\subsection{Tokenizace}
První krok kompileru je tokenizace -- vstupní soubor je rozdělen na jednotlivé tokeny,
jsou to klíčová slova (\code{while}, \code{if}), čísla (\code{15}, \code{1.79e308}), textové řetězce (\code{"Hello world!"}),
názvy funkcí, proměnných a~datových typů nebo symboly (\code{+}, \code{*}, \code{>=}, \code{;}).
V~tomto kroku se také zahodí komentáře (řádky začínající \code{//} nebo \code{/* komentar */}) a~všechny bílé znaky (mezery, tabulátory, konce řádků).

Ke každému tokenu se také přidá informace o~tom, kde se ve zdrojovém kódu vyskytuje. Pokud program narazí na chybu, může tak k~chybové hlášce přidat místo, kde se chyba vyskytla.

Pro tokenizaci jsem použil knihovnu Flex, která z~pravidel pro tokenizaci (\file{src/parse/lex.y}) vygeneruje potřebný kód.


\subsection{Syntaktická analýza}
Po tokenizaci přichází na řadu syntaktická analýza (neboli \uv{parsing}), která tokeny poskládá do takzvaného syntaktického stromu.
Každý vrchol v~syntaktickém stromě představuje nějakou operaci a~má pod sebou parametry této operace
(například operace \code{while} cyklu bude pod sebou mít podmínku a~kód cyklu).

\TODO{tikzni tu nějaký hezký obrázek kódu a jeho syntaktického stromu}

Syntaktický strom je v~mém kódu implementován ve složce \file{src/ast/};
v~souboru \file{src/ast/main.h} jsou definované datové typy pro jednotlivé vrcholy syntaktického stromu (např. \code{ast\_expr\_t} pro výrazy).
Celý strom je pak uložen v~datovém typu \code{ast\_t}. (Zkratka \code{ast} od anglického \uv{abstract syntax tree}.)

V~\file{src/ast/main.c} jsou implementace pomocných funkcí pro vytváření syntaktického stromu
a~v~\file{src/ast/print.c}, \file{src/ast/print.h} jsou funkce pro vypisování syntaktického stromu, slouží pouze pro debugovací účely. \TODO{smazat soubory?}

V~souborech \file{src/ast/scope.c} a~\file{src/ast/scope.h} je definovaný datový typ \code{ast\_scope\_t*}, který mapuje jména lokálních proměnných na jejich objekty.
Navíc si drží odkaz na vyšší \code{ast\_scope\_t*}, protože lokální prostory mohou být vložené do sebe (např. proměnnou definovanou uvnitř \code{for} cyklu nelze použít mimo tento cyklus).

K~samotnému převodu tokenů do syntaktického stromu jsem použil knihovnu GNU Bison.
Ta ze seznamu pravidel (\file{src/parse/parse.y}), které popisují jak se tokeny a~symboly mohou spojovat do jiných symbolů, vygeneruje kód, který tato pravidla implementuje.
Zárověn zkontroluje, že syntax jazyka je jednoznačný, tedy že každý seznam tokenů lze převést do syntaktického stromu pouze jedním způsobem.
To je důležité hlavně u~vyhodnocování výrazů, protože \code{a+b*c} se musí vyhodnotit jako \code{a+(b*c)}. Operace násobení, dělení a~modulo se musí vyhodnotit zleva doprava. \TODO{vic rozepsat, jake jsem toho dosahl, prip z toho udelat celou podkapitolu?}

Hlavní funkce \code{parse\_file} a~několik pomocných funkcí pro parsování jsou v~\file{src/parse/main.c}.

\subsection{Překlad do LLVM}
Ze syntaktického stromu se dále vygeneruje LLVM kód. Tento jazyk je něco mezi programovacím jazykem a~strojovým kódem.
Každá funkce v~tomto jazyce je seznamem takzvaných základních bloků, které jsou seznamem instrukcí.
Poslední instrukce v~bloku musí být skok do jiného bloku (často s~větvením podle hodnoty registru), návrat z~funkce nebo speciální \code{unreachable}, které definuje, že se program při spuštění nikdy
nedostane do tohoto stavu (toto umožňuje další optimalizace kódu).

Jednotlivé instrukce v~LLVM pracují s~registry.
Každá instrukce, která vrací nějakou hodnotu (tedy skoro všechny instrukce kromě ukládáí dat do paměti) rovnou zapíše tuto hodnotu do \textit{nového} registru.
Do registrů lze v~LLVM totiž zapisovat pouze jednou (to nám generování LLVM trochu ztěžuje, ale tato forma kódu je jednodušeji optimalizovatelnější).

Také lze použít pouze ty registry, které jsou \uv{přímými předchůdci} danému bloku kódu; to znamená, že pokud definujeme registr uvnitř \code{if} bloku, tak ho nemůžeme přímo použít mimo něj.
K~tomuto účelu slouží speciální instrukce \code{phi}, která sama o~sobě nic nepočítá, ale vrátí hodnotu podle toho, ze kterého bloku jsme se dostali do tohoto bloku.
Například funkce, která vrací větší ze dvou čísel by mohla vypadat takto:

\TODO{diakritika v code listingu}
\begin{lstlisting}[language=llvm]
i32 @max(i32 %x, i32 %y){
	entry:
		%cmp = icmp sgt i32 %x, %y  ; %cmp = 1 pokud %x < %y, jinak 0
		br i1 %cmp, label X, label Y
	X:
		br label end
	Y:
		br label end
	end:
		%res = phi i32 [%x, X], [%y, Y]  ; vrati %x pokud predchazejici blok byl X, jinak vrati %y
		ret i32 %res
}
\end{lstlisting}

Datové typy pro LLVM kód (bloky, funkce, registry) jsem si zadefinoval v~\file{src/llvm.h} a~funkce pro převod tohoto kódu do textové podoby jsou v~\file{src/llvm.c}.

Samotný převod ze syntaktického stromu do LLVM kódu jsem implementoval rekurzivně v~\file{src/ast2llvm.c}.
Při tomto převodu se také zachytí většina chyb, například operace se špatnými typy (sčítání čísla a~typu \code{void} apod.) nebo volání funkcí se špatným počtem nebo typem argumentů.

Původně jsem tento převod naprogramoval tak, že jsem hodnoty lokálních proměnných ukládal přímo do LLVM registrů a~pro smyčky a~podmínky jsem použil mnoho \code{phi} instrukcí.
To sice fungovalo bez problému, ale nebylo možné získat adresu proměnné v~paměti, jelikož proměnné nebyli uložené přímo v~paměti, ale v~registrech.
Ale protože jsem chtěl ve svém jazyku možnost získat adresy lokálních proměnných, musel jsem kód přepsat.
Nyní používám LLVM instrukci \code{alloca} abych na stacku rezervoval paměť pro proměnné, a~v~registrech ukládám pouze adresy na tyto proměnné.
Tím jsem ze zbavil i~skoro všech \code{phi} instrukcí, které však nahradilo množství instrukcí \code{store} (ukládání do paměti) a~\code{load} (čtení z~paměti).

\subsection{Kompilace LLVM kódu do strojového kódu}
Pro tuto část jsem použil již hotové programy z~LLVM toolchain a~linuxový linker \code{ld}.

Prvně se vygenerovaný LLVM kód prožene programem \file{opt}, který ho zoptimalizuje. Poté se optimalizovaný LLVM kód převede na strojový pomocí programu \file{llc}.
Do výsledného strojového kódu se pak pomocí \file{ld} přidá C~runtime library (ta před spuštěním samotného kódu připraví stack a~poté zavolá funkci \code{main} s~argumenty z~příkazové řádky)
a~\file{libc} (knihovna, která obsahuje moho užitečných funkcí -- bez ní by programy v~mém jazyce neuměli nic užitečného, protože můj jazyk neumí komunikovat s~operačním systémem).

\section{PROGRAMOVACÍ PROSTŘEDKY}
\TODO{JAKÝMI PROSTŘEDKY}\\
\textit{vývojové prostředí, použité knihovny, hotové programy atd. včetně zdůvodnění jejich výběru, použité při tvorbě maturitní práce}


\section{ZHODNOCENÍ DOSAŽENÝCH VÝSLEDKŮ}
\TODO{JAK SE TO POVEDLO}\\
\textit{výčet splněných a~nesplněných bodů zadání resp. vlastních předsevzetí}


\section{INSTALACE}
\TODO{JAK SE TO INSTALUJE}\\
\textit{popis instalace, systémové požadavky, nárok  na HW, SW, atp.}


\section{OVLÁDÁNÍ}
\subsection{Použití kompileru}

Kompiler se musí spustit z~příkazové řádky s~alespoň jedním argumentem: lokací vstupního souboru. (Kompiler umí kompilovat pouze jeden soubor najednou).
(Např.: \code{./compiler code.txt}.) V~tomto případě se kód zkompiluje, příp. vyhodí chybové hlášky nebo varování a~vypíše \uv{\code{no output path specified}}.

Pokud chceme výsledek kompilace uložit do souboru, stačí použít \uv{\code{-o <výstupni soubor>}}.

Kompiler umí generovat i~assembly nebo LLVM kód, pokud dostane vlajku \uv{\code{-filetype=asm}}, respektive \uv{\code{-filetype=llvm}}.
V~obou případech výpíše daný kód na standartní výstup, pokud nedostane lokaci výstupního souboru (vlajkou \uv{\code{-o}}).

Stupeň optimalizace lze ovládat vlajkami \code{-O0} (žádná optimalizace), \code{-O1}, \code{-O2} a~\code{-O3} (nejvyšší stupeň optimalizace).

\subsection{Popis jazyka}
Jazyk je velmi podobný jazyku C.
Kód se na nejvyšší úrovni skládá ze seznamu funkcí, globálních proměnných a~definicemi typů.

\subsubsection{Typy}
Jazyk podporuje pět hlavních typy typů: celá čísla, destinná čísla, pointery (adresy do paměti), strukturovaný datový typ a~speciální typ \code{void}.

Celočíselné typy jsou buď se znaménkem (\code{i8, i16, i32, i64}) nebo bez znaménka (\code{u8, u16, u32, u64}). Číslo v~jejich názvu odpovídá jejich velikosti (počtu bitů).
Speciální Celočíselné datový typ je navíc \code{bool}, který je vždy buď 1, nebo 0.

Destinné typy jsou dva: \code{f32} a~\code{f64}, které obsahují \code{float}, respektive \code{double}.

Typ pointer lze vytvořit z~libovolného typu přidáním \code{*} na konec. Například typ \code{float*} bude obsahovat adresu, na které se vyskytuje \code{float}.

Speciální typ \code{void} lze použít pouze při definici funkce, která nevrací žádnou hodnotu, nebo jako obecný pointer \code{void*}.

Strukturované datové typy lze definovat pomocí klíčového slova \code{struct}, složených závorek a~seznamem typů a~jmen, které má obsahovat.
Např.: \code{struct {i32 x; i32 y;}}.
Obsah proměnné se strukturovaným typem lze získat, stejně jako v~jazyku C, pomocí \uv\code{.}, nebo, pokud se jedná pointer k~strukturovanému typu, \uv\code{->}.

Typům lze také dát jiné jméno, pomocí klíčového slova \code{typedef}, následovaný typem, nové jméno pro daný typ a~středníkem. Toto se hodí hlavně u~strukturovaných typů.
Např.:
\begin{lstlisting}[language=C]
typedef i32 int;

typedef struct {
	int x;
	int y;
} vector2;
\end{lstlisting}

\subsubsection{Proměnné}
Lokální i~globální proměnné se definují typem proměnné, jménem, případně počáteční hodnotou (pouze u~lokálních proměnných) a~středníkemů.
Např.: \code{u64 number;} nebo \code{u32 x = y*5};

U~globalních proměnných nelze nastavit počáteční hodnotu. \TODO{odeber to z parseru a vlastne odevsad}

Adresu proměnné lze získat operací \code{\&}. (Např.: \code{u8* ptr = \&number;})

\subsubsection{Funkce}
Funkci lze buď definovat, nebo deklarovat. Deklarace pouze řekne kompileru, jaké argumenty funkce očekává a~jaký datový typ vrací.
Jsou potřeba při použití funkcí z~externí knihovny nebo rekurzivních funkcí (ve funkci lze použít pouze ty funkce a~globální proměnné, které jsou deklarované nebo definované dříve v~kódu).

Deklarace funkce začíná typem, který funkce vrací (nebo speciální typ \code{void}, pokud funce nevrací žádnou hodnotu),
následuje jméno funkce, pak v~kulatých závorkách typy a~jména argumentů a~nakonec středník.

Například:
\begin{lstlisting}[language=C]
i32 puts(u8* s);
void* malloc(u64 size);
void free(void* ptr);
\end{lstlisting}

Definice funkce začíná stejně jako deklarace funkce, ale místo středníku následuje tělo funkce ve složených závorkách.
Například:

\begin{lstlisting}[language=C]
i32 max(i32 x, i32 y){
	if (x > y){
		return x;
	} else {
		return y;
	}
}
\end{lstlisting}

\subsubsection{Podmínky}
Podmínky se definuje pomocí klíčových slov \code{if} a~\code{else}.
Po \code{if} vždy následuje výraz v~kulatých závorkách podle kterého se program rozhodne, kterou část podmínky vyhodnotí.
Poté se vyskutje tělo \code{if}u, případně \code{else} a~další tělo. \TODO{je "telo" spravne slovo pro tohle?}

V~obou případech je tělo buď jeden řádek kódu (končící prvním středníkem), nebo libovolný počet řádků obalených ve složených závorkách.

\subsubsection{Výrazy}
Výrazem se myslí jakákoliv proměnná, číselná konstanta nebo složení dvou výrazů nějakou operací.

Povolené operace jsou:
\begin{itemize}
	\item násobení, děleno a~modulo (\code{*}, \code{/}, \code{\%})
	\item sčítání a~odčítání (\code{+}, \code{-})
	\item bitový posun doleva a~doprava (\code{<<}, \code{>>})
	\item porovnání (\code{==}, \code{!=}, \code{<}, \code{>}, \code{<=}, \code{>=})
	\item bitwise AND (\code{\&})
	\item bitwise XOR (\code{\^})
	\item bitwise OR  (\code{|})
	\item logický AND (\code{\&\&}) -- druhý argument se vyhodnotí pouze pokud se první vyhodnotí jako 1
	\item logický OR (\code{||}) -- druhý argument se vyhodnotí pouze pokud se první vyhodnotí jako 0
\end{itemize}
Operace se vyhodnocují v~pořadí tohoto seznamu zeshora (tedy \code{a || b + c | d} se vyhodnotí jako \code{a || ((b+c) | d)}).
Operace ve stejné skupině se vyhodnocují vždy zleva doprava.

Navíc jsou povolené operace s~jedním argumentem:
\begin{itemize}
	\item negace (\code{-})
	\item bitwise NOT (\code{\~})
	\item logický NOT (\code{!})
	\item dereference (\code{*}) -- čtení z~paměti, funguje pouze s~typem pointer
\end{itemize}

Výrazy lze libovolně závorkovat.


\section{SEZNAM POUŽITÝCH INFORMAČNÍCH ZDROJŮ}
...


\showtodo
\end{document}
